%!TeX root=csp.tex
\section{Vocabulary}\label{sec:vocabulary}

This section fixes the universal-algebraic vocabulary. It is Zhuk's
\S2.1, rewritten so that every definition is unambiguous when read by a
machine. The substantive content begins in \S\ref{sec:types}; a reader who
knows the subject should skim to
Definition~\ref{def:irreducible} (irreducible congruences) and
Definition~\ref{def:bridge} (bridges), which are the two places where the
conventions genuinely matter.

\subsection{Algebras, subuniverses, terms}

\begin{definition}[Algebra]\label{def:algebra}
An algebra is a pair $\mathbf A=(A;w^{\mathbf A})$ with $A$ a finite nonempty
set and $w^{\mathbf A}\colon A^{n}\to A$, subject to
Convention~\ref{conv:standing}. We write $A$ for the universe of $\mathbf A$
and drop the superscript on $w$ when no confusion arises.
\end{definition}

\begin{definition}[The algebras $\mathbf Z_{p}$]\label{def:Zp}
For a prime $p$ with $p\mid n-1$, $\mathbf Z_{p}$ is the algebra on
$\{0,1,\dots,p-1\}$ with $w^{\mathbf Z_{p}}(x_{1},\dots,x_{n})
=x_{1}+\dots+x_{n}\bmod p$.
\end{definition}

\begin{hazard}[$\mathbf Z_{p}$ is only sometimes in $\Vn$]\label{haz:Zp}
Zhuk defines $\mathbf Z_{p}$ for every prime $p$ and then asserts that
``every algebra $\mathbf Z_{p}$ belongs to $\Vn$''. That assertion carries a
silent side condition. The operation $x_{1}+\dots+x_{n}\bmod p$ is idempotent
exactly when $n\equiv 1\pmod p$: idempotence says $nx\equiv x$ for all $x$,
i.e.\ $p\mid n-1$. Given that, it is also a WNU (all the identities read
$(n-1)x+y=y$) and special (both sides reduce to $y$).

So $\mathbf Z_{p}\in\Vn$ if and only if $p\mid n-1$, and every statement of the
form ``$\mathbf A/\sigma\cong\mathbf Z_{p}$ for some prime $p$''
(Lemma~\ref{lem:linear-on-top}, Definition~\ref{def:perfect-linear}) silently
produces a $p$ dividing $n-1$. A formalization must decide whether
$\mathbf Z_{p}$ is defined for all $p$ and merely fails to lie in $\Vn$ for bad
$p$, or is defined only for $p\mid n-1$. The first is cleaner; the second makes
the membership free. Either way the side condition must be visible, because
``for some prime $p$'' otherwise reads as an unconstrained existential.
\end{hazard}

\begin{definition}[Subuniverse, subalgebra]\label{def:subuniverse}
$B\subseteq A$ is a \emph{subuniverse} of $\mathbf A$ if
$w(b_{1},\dots,b_{n})\in B$ for all $b_{1},\dots,b_{n}\in B$. If moreover
$B\neq\varnothing$ then $\mathbf B=(B;w\!\restriction_{B^{n}})$ is a
\emph{subalgebra}, written $\mathbf B\le\mathbf A$.
\end{definition}

\begin{definition}[Term operations]\label{def:clone}
$\Clo(\mathbf A)$ is the least set of finitary operations on $A$ containing all
projections and closed under composition with $w$. Its $m$-ary part is written
$\Clo_{m}(\mathbf A)$. Elements of $\Clo(\mathbf A)$ are the \emph{term
operations} of $\mathbf A$.
\end{definition}

\begin{hazard}[Terms versus term operations]\label{haz:term-vs-op}
Zhuk writes $t\in\Clo(\mathbf A)$ and treats $t$ as both a syntactic term and
the operation it induces. These must be distinguished in a formalization: a
term is a finite tree, and distinct terms may induce the same operation. Every
statement below that quantifies over $\Clo(\mathbf A)$ is a statement about
\emph{operations}, so the syntactic layer may be quotiented away --- but it
cannot be skipped, because $\Sg$ is characterised through terms
(Lemma~\ref{lem:sg-terms}) and because the arity bookkeeping in the absorption
arguments is syntactic.
\end{hazard}

\begin{definition}[Generation]\label{def:sg}
For $X\subseteq A^{m}$, $\Sg_{\mathbf A}(X)$ is the least subuniverse of
$\mathbf A^{m}$ containing $X$.
\end{definition}

\begin{lemma}[Generation by terms]\label{lem:sg-terms}
Let $X=\{x_{1},\dots,x_{k}\}\subseteq A^{m}$ be nonempty. Then
\[
  \Sg_{\mathbf A}(X)=\bigl\{\,t^{\mathbf A^{m}}(x_{1},\dots,x_{k})
      \;:\;t\in\Clo_{k}(\mathbf A)\,\bigr\}.
\]
The generator list is \emph{fixed}: the same list $x_{1},\dots,x_{k}$ in the
same order serves every element of the closure, with $t$ varying.
\end{lemma}

\begin{hazard}\label{haz:sg-fixed-list}
The ``fixed list'' phrasing is what makes Lemma~\ref{lem:sg-terms} usable. The
weaker reading --- for each $y$ in the closure there are \emph{some}
generators and \emph{some} term --- is also true but useless, because every
later application applies one term simultaneously to several closure elements
and needs their generator lists to coincide. This was a defect found in review
of the predecessor blueprint and is called out here for the same reason.
\end{hazard}

\subsection{Relations}

\begin{definition}[Relational vocabulary]\label{def:rel-vocab}
Let $R\subseteq A_{1}\times\dots\times A_{m}$.
\begin{enumerate}\alphlabels\tightlist
\item $R$ is \emph{subdirect} if $\proj_{i}(R)=A_{i}$ for every $i$;
  $\mathbf R\sdle\mathbf A_{1}\times\dots\times\mathbf A_{m}$ abbreviates
  ``$R$ is a subdirect subuniverse''.
\item For $R\subseteq A^{m}$: $R$ is \emph{reflexive} if
  $(a,\dots,a)\in R$ for every $a\in A$.
\item For binary $\delta_{1}\subseteq A_{1}\times A_{2}$,
  $\delta_{2}\subseteq A_{2}\times A_{3}$:
  $\delta_{1}\circ\delta_{2}=\{(a,c):\exists b\;(a,b)\in\delta_{1}\wedge
  (b,c)\in\delta_{2}\}$, and $\delta^{-1}=\{(b,a):(a,b)\in\delta\}$.
\item For $B\subseteq A_{1}$ and $\delta\subseteq A_{1}\times A_{2}$:
  $B\circ\delta=\{c:\exists b\in B\;(b,c)\in\delta\}$.
\item A subdirect $\delta\subseteq A\times B$ is \emph{linked} if the bipartite
  graph on $A\sqcup B$ with edge set $\delta$ is connected. It is
  \emph{bijective} if $|\delta|=|A|=|B|$.
\item For binary $\sigma$ and $m\ge 2$,
  $\sigma^{[m]}=\{(a_{1},\dots,a_{m}):(a_{i},a_{j})\in\sigma\ \forall i,j\}$.
\end{enumerate}
\end{definition}

\begin{hazard}[Composition is diagrammatic]\label{haz:comp-order}
$\delta_{1}\circ\delta_{2}$ applies $\delta_{1}$ \emph{first}. This is the
order of \texttt{Rel.comp} and the opposite of \texttt{Function.comp}. Every
composite of bridges and of congruences below is written left to right, and
silently reversing it produces statements that are true-looking and wrong ---
in Corollary~\ref{cor:stable-intersection}(l) the composite
$\sigma_{k}\circ\proj_{k,\ell}(R)\circ\sigma_{\ell}$ is asymmetric in $k$ and
$\ell$ and the order is the content.
\end{hazard}

\begin{hazard}[``Linked'' is used for two different things]\label{haz:linked}
Zhuk uses \emph{linked} both for a binary relation (Definition~\ref{def:rel-vocab}(e))
and for a CSP instance (Definition~\ref{def:instance-linked}). They are related
but not the same predicate, and the coincidence of names causes real confusion
in \S\ref{sec:main-statements}, where both occur in one proof. Use two names.
\end{hazard}

\begin{definition}[Quotients of relations]\label{def:quotient-rel}
For an equivalence $\sigma$ on $A$ and $a\in A$, $a/\sigma$ is the class of
$a$; for $B\subseteq A$, $B/\sigma=\{b/\sigma:b\in B\}$; for
$R\subseteq A^{m}$, $R/\sigma=\{(b_{1}/\sigma,\dots,b_{m}/\sigma):\bar b\in R\}$.
\end{definition}

\begin{hazard}[Images do not commute with intersection]\label{haz:quotient-rel}
$B/\sigma$ and $R/\sigma$ are \emph{images}. Three consequences that the source
navigates silently:
\begin{enumerate}\alphlabels\tightlist
\item $B/\sigma$ is a subset of $A/\sigma$, not the quotient of $B$ by
  $\sigma\cap B^{2}$. For $B$ a subuniverse and $\sigma$ a congruence the two
  are canonically isomorphic as algebras, and the source moves between them
  freely. Fix $B/\sigma$ to be the image and prove the isomorphism once.
\item In general $(C_{1}\cap C_{2})/\sigma\subsetneq C_{1}/\sigma\cap
  C_{2}/\sigma$. Lemma~\ref{lem:propagation}(fm) needs equality for a family
  $C_{1},\dots,C_{t}$, and gets it from the extra hypothesis
  $(C_{i}\circ\sigma)\cap B=C_{i}$, i.e.\ that each $C_{i}$ is
  $\sigma$-saturated inside $B$. That is a content step, not notation.
\item $R/\sigma$ need not be stable under anything, and need not be a
  subuniverse of $(\mathbf A/\sigma)^{m}$ unless $\sigma$ is a congruence.
\end{enumerate}
\end{hazard}

\subsection{Rectangularity}

\begin{definition}[Parallelogram property, rectangularity]\label{def:parallelogram}
Let $R\subseteq A_{1}\times\dots\times A_{m}$.
\begin{enumerate}\alphlabels\tightlist
\item The $i$-th variable of $R$ is \emph{rectangular} if for all
  $\bar a,\bar b$,
  \[
    \bar a[i\mapsto b_{i}]\in R,\quad
    \bar b[i\mapsto a_{i}]\in R,\quad
    \bar b\in R
    \;\Longrightarrow\;\bar a\in R,
  \]
  where $\bar a[i\mapsto c]$ is $\bar a$ with its $i$-th entry replaced by $c$.
  $R$ is \emph{rectangular} if every variable is.
\item $R$ \emph{has the parallelogram property} if for every permutation of the
  coordinates giving $R'$ and every $\ell\in[m-1]$,
  \[
    (a_{1},\dots,a_{\ell},b_{\ell+1},\dots,b_{m})\in R',\quad
    (b_{1},\dots,b_{\ell},a_{\ell+1},\dots,a_{m})\in R',\quad
    \bar b\in R'
    \;\Longrightarrow\;\bar a\in R'.
  \]
\item The \emph{rectangular closure} of $R$ is the least rectangular relation
  containing $R$.
\end{enumerate}
\end{definition}

\begin{lemma}[Two assertions the source makes in passing]\label{lem:rect-basics}
\begin{enumerate}\alphlabels\tightlist
\item A relation with the parallelogram property is rectangular.
\item The rectangular closure exists.
\end{enumerate}
\end{lemma}

\begin{proof}
(a) Apply the parallelogram property to the permutation carrying coordinate
$i$ to position $1$, with $\ell=1$. (b) Rectangularity at coordinate $i$ is a
Horn implication, so the rectangular relations containing $R$ are closed under
arbitrary intersection and the family is nonempty (it contains
$A_{1}\times\dots\times A_{m}$); the intersection of all of them is the least
one.
\end{proof}

\begin{hazard}\label{haz:rect-closure}
Both parts are asserted without proof in \cite{ZhukSimplified} --- (a) as ``as
it follows from the definitions'', (b) by the definite article in ``\emph{the}
minimal rectangular relation''. Both are easy and both must be present. Note
that ``$R$ is rectangular'' unqualified --- the form used in ``rectangular
closure'' --- is never defined in the source; it means every variable is
rectangular.
\end{hazard}

\begin{hazard}[Quantifying over permutations]\label{haz:permutations}
The parallelogram property is stated ``any permutation of its variables gives a
relation $R'$ satisfying \dots''. Formally this is a quantification over
$\mathfrak S_{m}$ combined with a quantification over the split point $\ell$;
equivalently, and more usably, over all ways of splitting $[m]$ into two
nonempty blocks. The second formulation is the one to formalize: it avoids
transporting $R$ along a permutation, which in a dependently typed setting
means an equivalence of index types and a coercion at every use.
\end{hazard}

\subsection{Congruences}

\begin{definition}[Stability]\label{def:stable}
Let $R\subseteq A_{1}\times\dots\times A_{m}$ and let $\sigma$ be a binary
relation on $A_{i}$. The $i$-th variable of $R$ is \emph{stable under $\sigma$}
if $\bar a\in R$ and $(a_{i},b)\in\sigma$ imply $\bar a[i\mapsto b]\in R$.
$R$ is \emph{stable under $\sigma$} if every variable is (each with respect to
$\sigma$ on the corresponding factor, when the factors agree).
\end{definition}

\begin{definition}[Congruence]\label{def:congruence}
A \emph{congruence} on $\mathbf A$ is an equivalence relation on $A$ that is a
subuniverse of $\mathbf A\times\mathbf A$. It is \emph{nontrivial} if it is
neither the equality relation $0_{\mathbf A}$ nor $A^{2}$. The quotient
$\mathbf A/\sigma$ is the algebra on $A/\sigma$ with
$w(\bar a/\sigma)=w(\bar a)/\sigma$.
\end{definition}

\begin{definition}[Irreducible congruence and its cover]\label{def:irreducible}
A congruence $\sigma$ on $\mathbf A$ is \emph{irreducible} if it cannot be
written as an intersection of binary subuniverses of $\mathbf A\times\mathbf A$
each stable under $\sigma$ and each properly containing $\sigma$. Equivalently:
there are no $\mathbf S_{1},\dots,\mathbf S_{k}\le\mathbf A/\sigma\times
\mathbf A/\sigma$ with $0_{\mathbf A/\sigma}=S_{1}\cap\dots\cap S_{k}$ and
$S_{i}\neq 0_{\mathbf A/\sigma}$ for every $i$.

For irreducible $\sigma$, $\cover{\sigma}$ denotes the least
$\delta\le\mathbf A\times\mathbf A$ with $\delta\supsetneq\sigma$ and $\delta$
stable under $\sigma$.
\end{definition}

\begin{hazard}[Three traps in Definition~\ref{def:irreducible}]\label{haz:irreducible}
\begin{enumerate}\alphlabels\tightlist
\item Irreducibility is \emph{not} meet-irreducibility in the congruence
  lattice. The $\mathbf S_{i}$ range over binary \emph{subuniverses} stable
  under $\sigma$, which need not be equivalence relations. The two notions
  differ, and every use below is the stated one.
\item $\cover{\sigma}$ need not be a congruence. It is a subuniverse of
  $\mathbf A^{2}$ stable under $\sigma$, nothing more; that $\cover{\sigma}$
  \emph{is} a congruence is an extra hypothesis in the definition of a linear
  congruence (Definition~\ref{def:linear-congruence}), which would be
  redundant otherwise.
\item Existence and uniqueness of $\cover{\sigma}$ need an argument, given in
  Proposition~\ref{prop:cover}. Uniqueness is exactly what irreducibility buys;
  without it there is no ``the'' cover. A formalization must produce
  $\cover{\sigma}$ as data attached to a proof of irreducibility, not as a
  standalone function.
\item The family $S_{1},\dots,S_{k}$ may be \emph{empty}
  (item \ref{leg:emptyfamily} of
  \S\ref{sec:legislation}). This is not a free choice: the two formulations
  above agree only under it.
\end{enumerate}
\end{hazard}

\begin{proposition}[The cover exists, and is a tolerance]\label{prop:cover}
Let $\sigma$ be an irreducible congruence on $\mathbf A$. Then
\begin{enumerate}\alphlabels\tightlist
\item $\sigma\neq A^{2}$;
\item the family $\mathcal F$ of $\delta\le\mathbf A\times\mathbf A$ with
  $\delta\supsetneq\sigma$ and $\delta$ stable under $\sigma$ is nonempty and
  closed under intersection, so it has a least element $\cover{\sigma}$;
\item $\cover{\sigma}$ is reflexive and symmetric.
\end{enumerate}
\end{proposition}

\begin{proof}
(a) The empty family is allowed, and its intersection is $A^{2}$; so if
$\sigma=A^{2}$ then $\sigma$ is an intersection of a family of relations each
properly containing it, vacuously, and $\sigma$ is reducible.

(b) $\mathcal F\ni A^{2}$ by (a). If $\delta_{1},\delta_{2}\in\mathcal F$ then
$\delta_{1}\cap\delta_{2}$ is again a subuniverse of $\mathbf A^{2}$ stable
under $\sigma$ and contains $\sigma$; it contains $\sigma$ \emph{properly},
since otherwise $\sigma=\delta_{1}\cap\delta_{2}$ exhibits $\sigma$ as an
intersection of two members of $\mathcal F$, contradicting irreducibility. So
$\mathcal F$ is closed under binary, hence finite, intersection, and
$\cover{\sigma}=\bigcap\mathcal F\in\mathcal F$ is its least element.

(c) $\cover{\sigma}\supseteq\sigma\supseteq 0_{\mathbf A}$ is reflexive. For
symmetry, $(\cover{\sigma})^{-1}$ is a subuniverse of $\mathbf A^{2}$, is
stable under $\sigma$ because $\sigma$ is symmetric, and properly contains
$\sigma^{-1}=\sigma$; so $(\cover{\sigma})^{-1}\in\mathcal F$ and therefore
$\cover{\sigma}\subseteq(\cover{\sigma})^{-1}$, whence equality.
\end{proof}

\begin{remark}[Why not a congruence]\label{rem:cover-not-cong}
Nothing in Proposition~\ref{prop:cover} gives transitivity, and the source
supplies the internal evidence that none is available: that $\cover{\sigma}$ is
a congruence is item (b) of Definition~\ref{def:linear-congruence} and
conclusion (a) of Lemma~\ref{lem:nontrivial-bridge}, both of which would be
redundant otherwise. Symmetry of $\cover{\sigma}$, by contrast, is used
silently in \cite{ZhukSimplified} --- in the proof of
Lemma~\ref{lem:pc-bridges-trivial}, where the argument is run again ``switching
$x_{1}$ and $x_{2}$''.
\end{remark}

\begin{lemma}[Passing to the quotient]\label{lem:irred-quotient}
Let $\sigma$ be a congruence on $\mathbf A$. Then $\delta\mapsto\delta/\sigma$
is an inclusion-preserving bijection, commuting with intersections, from the
subuniverses of $\mathbf A^{2}$ that contain $\sigma$ and are stable under
$\sigma$ onto the subuniverses of $(\mathbf A/\sigma)^{2}$, sending $\sigma$ to
$0_{\mathbf A/\sigma}$. Consequently $\sigma$ is irreducible, linear or PC
exactly when $0_{\mathbf A/\sigma}$ is, and then
$\cover{\sigma}/\sigma=\cover{(0_{\mathbf A/\sigma})}$.
\end{lemma}

\begin{proof}
A subuniverse $\delta$ of $\mathbf A^{2}$ containing $\sigma$ and stable under
$\sigma$ is a union of products of $\sigma$-blocks, so it is the preimage of
$\delta/\sigma$; conversely the preimage of a subuniverse of
$(\mathbf A/\sigma)^{2}$ is a subuniverse containing $\sigma$ and stable under
$\sigma$. The two constructions are mutually inverse and preserve inclusion and
intersection, and $\sigma$ is the preimage of $0_{\mathbf A/\sigma}$. The
statement of Definition~\ref{def:irreducible} is a statement about that family
alone, so it transports; and both the least element in
Proposition~\ref{prop:cover}(b) and the data of
Definition~\ref{def:linear-congruence} transport with it.
\end{proof}

\begin{hazard}[This is the licence for ``assume $\sigma=0$'']\label{haz:factor-out}
Three proofs in \S\ref{sec:properties} open with ``we may factor by $\sigma$
and assume $\sigma$ is the equality relation''.
Lemma~\ref{lem:irred-quotient} is what licenses that, and it is one of the
sentences \cite{ZhukSimplified} states without proof
(\S\ref{sec:legislation}, item \ref{leg:notice}).
\end{hazard}

\subsection{Bridges}

\begin{definition}[Bridge]\label{def:bridge}
Let $\sigma_{1},\sigma_{2}$ be congruences on $\mathbf D_{1},\mathbf D_{2}$. A
relation $\delta\le\mathbf D_{1}^{2}\times\mathbf D_{2}^{2}$ is a \emph{bridge
from $\sigma_{1}$ to $\sigma_{2}$} if
\begin{enumerate}\alphlabels\tightlist
\item the first two variables of $\delta$ are stable under $\sigma_{1}$;
\item the last two variables of $\delta$ are stable under $\sigma_{2}$;
\item $\proj_{1,2}(\delta)\supsetneq\sigma_{1}$ and
      $\proj_{3,4}(\delta)\supsetneq\sigma_{2}$;
\item $(a_{1},a_{2},a_{3},a_{4})\in\delta$ implies
      $\bigl((a_{1},a_{2})\in\sigma_{1}\iff(a_{3},a_{4})\in\sigma_{2}\bigr)$.
\end{enumerate}
For a bridge $\delta$ put $\bcol{\delta}=\{(x,y):(x,x,y,y)\in\delta\}$.
A bridge $\delta\subseteq D^{4}$ is \emph{reflexive} if $(a,a,a,a)\in\delta$
for all $a\in D$.
\end{definition}

\begin{remark}[The motivating example]\label{rem:bridge-example}
On $\mathbb Z_{4}$, the relation
$\delta=\{(a_{1},a_{2},a_{3},a_{4}):a_{1}-a_{2}=2a_{3}-2a_{4}\}$ is a bridge
from $0_{\mathbb Z_{4}}$ to congruence mod~$2$. Condition (d) is the content:
$a_{1}=a_{2}$ exactly when $a_{3}\equiv a_{4}\pmod 2$. Bridges are the
formal shadow of a centralizer relation; see \cite{RossSlides}.
\end{remark}

\begin{definition}[Composition of bridges]\label{def:bridge-comp}
For bridges $\delta_{1}$ from $\sigma_{0}$ to $\sigma_{1}$ and $\delta_{2}$
from $\sigma_{1}$ to $\sigma_{2}$, set
\[
  (\delta_{1}\ast\delta_{2})(x_{1},x_{2},z_{1},z_{2})
  =\exists y_{1}\exists y_{2}\;
   \delta_{1}(x_{1},x_{2},y_{1},y_{2})\wedge\delta_{2}(y_{1},y_{2},z_{1},z_{2}).
\]
\end{definition}

\begin{lemma}[Bridge composition; {\cite[Lem.~6.3]{ZhukJACM}}]\label{lem:bridge-comp}
If $\sigma_{0},\sigma_{1},\sigma_{2}$ are irreducible then
$\delta_{1}\ast\delta_{2}$ is a bridge from $\sigma_{0}$ to $\sigma_{2}$, and
$\bcol{\delta_{1}\ast\delta_{2}}=\bcol{\delta_{1}}\circ\bcol{\delta_{2}}$.
\end{lemma}

\begin{hazard}\label{haz:bridge-comp}
Lemma~\ref{lem:bridge-comp} is imported, not proved, in
\cite{ZhukSimplified}; it is \cite[Lemma 6.3]{ZhukJACM}. Both conclusions are
used: the first in every chain of bridges, the second --- the identity
$\bcol{\cdot}$ commutes with composition --- in
Lemma~\ref{lem:perfect-from-linked}, where a bridge with $\bcol{\delta}$ linked
is powered up until $\bcol{\delta}=A^{2}$. The irreducibility hypothesis on
\emph{all three} congruences is needed and is easy to drop by accident.
\end{hazard}

\begin{definition}[Symmetric and pair-reflexive bridges]\label{def:pair-reflexive}
Let $\delta$ be a bridge from $\sigma$ to $\sigma$ on $\mathbf D$. Call $\delta$
\emph{symmetric} if
$\delta(x_{1},x_{2},x_{3},x_{4})\iff\delta(x_{3},x_{4},x_{1},x_{2})$, and
\emph{pair-reflexive} if
\[
  (x_{1},x_{2})\in\proj_{1,2}(\delta)\;\Longrightarrow\;
  (x_{1},x_{2},x_{1},x_{2})\in\delta .
\]
Write $\delta^{-1}(x_{1},x_{2},x_{3},x_{4})=\delta(x_{3},x_{4},x_{1},x_{2})$.
\end{definition}

\begin{lemma}[Symmetrisation]\label{lem:symmetrise}
Let $\delta$ be a reflexive bridge from $\sigma$ to $\sigma$ on $\mathbf D$ and
put $\delta^{\circ}=\delta\ast\delta^{-1}$, that is
\[
  \delta^{\circ}(x_{1},x_{2},x_{3},x_{4})=\exists x_{5}\exists x_{6}\;
  \delta(x_{1},x_{2},x_{5},x_{6})\wedge\delta(x_{3},x_{4},x_{5},x_{6}).
\]
Then $\delta^{\circ}$ is a reflexive, symmetric, pair-reflexive bridge from
$\sigma$ to $\sigma$ with
$\proj_{1,2}(\delta^{\circ})=\proj_{3,4}(\delta^{\circ})=\proj_{1,2}(\delta)$
and $\bcol{\delta^{\circ}}=\bcol{\delta}\circ\bcol{\delta}^{-1}
\supseteq\bcol{\delta}$. Moreover any pair-reflexive bridge $\varepsilon$
satisfies $\varepsilon\subseteq\varepsilon\ast\varepsilon$.
\end{lemma}

\begin{proof}
Symmetry is visible in the formula. For pair-reflexivity, if
$(x_{1},x_{2})\in\proj_{1,2}(\delta^{\circ})$ then
$\delta(x_{1},x_{2},x_{5},x_{6})$ for some $x_{5},x_{6}$, and taking the same
witnesses twice gives $(x_{1},x_{2},x_{1},x_{2})\in\delta^{\circ}$; the same
computation shows $\proj_{1,2}(\delta^{\circ})=\proj_{1,2}(\delta)$, the
inclusion $\subseteq$ being immediate, and symmetry then gives
$\proj_{3,4}(\delta^{\circ})=\proj_{1,2}(\delta^{\circ})$. Clause (c) of
Definition~\ref{def:bridge} follows, and clause (d) because
$(x_{1},x_{2})\in\sigma\iff(x_{5},x_{6})\in\sigma\iff(x_{3},x_{4})\in\sigma$.
Stability is inherited from $\delta$, and reflexivity from
$(a,a,a,a)\in\delta$. The identity
$\bcol{\delta^{\circ}}=\bcol{\delta}\circ\bcol{\delta^{-1}}
=\bcol{\delta}\circ\bcol{\delta}^{-1}$ is Lemma~\ref{lem:bridge-comp};
it contains $\bcol{\delta}$ because $\bcol{\delta}$ is reflexive, so
$(x,y)\in\bcol{\delta}$ gives $(x,y)\in\bcol{\delta}\circ\bcol{\delta}^{-1}$ by
taking $y$ as the intermediate point. Finally, if $\varepsilon$ is
pair-reflexive and $\varepsilon(x_{1},x_{2},x_{3},x_{4})$, then taking
$(x_{1},x_{2})$ itself as the intermediate pair --- legitimate because
$\varepsilon(x_{1},x_{2},x_{1},x_{2})$ --- gives
$(\varepsilon\ast\varepsilon)(x_{1},x_{2},x_{3},x_{4})$.
\end{proof}

\begin{hazard}[Pair-reflexivity is the hypothesis three lemmas need]\label{haz:pair-reflexive}
\cite{ZhukSimplified} states pair-reflexivity only once, as a hypothesis of
Lemma~\ref{lem:pc-inductive-step}, but three further statements require it and
do not say so; without it one of them is false. See
Hazard~\ref{haz:bridge-repairs}. Lemma~\ref{lem:symmetrise} is what makes this
costless: every bridge these statements are ever applied to is a
$\delta^{\circ}$.
\end{hazard}

\begin{definition}[Perfect linear congruence]\label{def:perfect-linear}
A congruence $\sigma$ on $\mathbf A$ is a \emph{perfect linear congruence} if
it is irreducible and there are a prime $p$ and
$\zeta\le\mathbf A\times\mathbf A\times\mathbf Z_{p}$ with
$\proj_{1,2}(\zeta)=\cover{\sigma}$ and
\[
  (a_{1},a_{2},b)\in\zeta\;\Longrightarrow\;
  \bigl((a_{1},a_{2})\in\sigma\iff b=0\bigr).
\]
\end{definition}

Perfect linear congruences are the payoff of the whole bridge machinery: they
let the passage from $\sigma$ to $\cover{\sigma}$ be tracked by a single
element of $\mathbf Z_{p}$, which is what turns a CSP instance into a system of
linear equations in \S\ref{sec:codim-one}.
